\scanpage{31}
\mspage{\handwritten{\unsure{63}}}
signaler en remarque). Par suite, $\operatorname{Ker}i$ est isomorphe à $N_{k(\underline p)/k}$, d'où
\[
\operatorname{rang}\operatorname{Im}i=r-t=n-s-t
\]
où
\[
\begin{aligned}
n&=\dim A,\quad \handwritten{s}=\dim A/\underline p,\quad r=n-s=\dim A_{\underline p}=\operatorname{rang}\underline p/\underline p^{2}\otimes k(t)\\
t&=\operatorname{rang}N_{k(p)/k}
\end{aligned}
\]
D'autre part, la deuxième ligne de $(D_{\handwritten{\alpha}})$ nous prouve que
\[
\operatorname{rang}\operatorname{Im}i_{\handwritten{\alpha}}
=\operatorname{rang}\handwritten{\hat\Omega}^{1}_{A/K_{\handwritten{\alpha}}}\otimes k(p)
-\operatorname{rang}\handwritten{\hat\Omega}^{1}_{k(\underline p)/K_{\handwritten{\alpha}}}
\]
Or \deleted{le premier terme du second membre, compte tenu que $A$ est \unsure{formellem|formelle}} sauf erreur il doit être vrai que pour $\handwritten{\alpha}$ grand, $\handwritten{\hat\Omega}^{1}_{A/K_{\handwritten{\alpha}}}$ est un module \handwritten{\uline}{libre de rang} $n+\operatorname{rang}\handwritten{\Omega}^{1}_{K_{\unsure{o|0}}/K_{\handwritten{\alpha}}}$ sur $A$. Il est vrai que je ne l'ai pas prouvé, et \deleted{\unsure{étant}} voulant en finir, je vais faire des hypothèses plus restrictives dans l'énoncé de 8.1. savoir que $A$ est formellement simple, non seulement sur $k$, mais même sur $K_o$ (il faut donc modifier la rédaction de l'énoncé 8.1. à moins que tu ne te sentes le courage de résoudre le point que je laisse en suspens ici). Dans ce cas, l'assertion que j'ai faite sur $\handwritten{\hat\Omega}^{1}_{A/K}$ est même vrai pour \uline{tout} sous-corps $K$ d'indice fini dans $K_o$, et se vérifie sans difficulté. \deleted{\unsure{E|K}}\unsure{C|K}omparant alors les deux expressions trouvées pour $\operatorname{rang}\operatorname{Im}i_{\handwritten{\alpha}}$ et $\operatorname{rang}\operatorname{Im}i$, on est remené à montrer que les deux sont égaux pour $\handwritten{\alpha}$ grand, i.e. qu'on a pour $\handwritten{\alpha}$ grand l'égalité
\[
\operatorname{rang}\handwritten{\hat\Omega}^{1}_{k(\underline p)/K_{\handwritten{\alpha}}}
-\operatorname{rang}\handwritten{\Omega}^{1}_{K_o/K_{\handwritten{\alpha}}}
=s+t=\dim A/\underline p+\operatorname{rang}N_{\handwritten{k(p)/k}}
\]
Pour abréger, posons $L=k(\underline p)$, \handwritten{no}tons que c'est une extension finie d'une extension de la forme $\handwritten{\bar}{K}_o=K_o((y_1,\ldots,y_s))$, et que l'on a
\[
\handwritten{\hat\Omega}^{1}_{L/K_{\handwritten{\alpha}}}=\handwritten{\Omega}^{1}_{L/K'_{\handwritten{\alpha}}}
\]
où on pose
\[
K'_{\handwritten{\alpha}}=K_{\handwritten{\alpha}}((y_1^p,\ldots,y_s^p)),
\]
ce qui ramène les modules avec chapeau à des modules sans. On est ainsi
% archive sheet number: 30
